\documentclass[11pt]{article}
\usepackage[utf8]{inputenc}
\usepackage[T1]{fontenc}
\usepackage{lmodern}
\usepackage[spanish]{babel}
\usepackage{amsmath,amssymb}
\usepackage{graphicx}
\usepackage[left=2cm,right=2cm,top=2cm,bottom=2cm]{geometry}
\usepackage{enumitem}
\usepackage{booktabs}
\usepackage{array}
\usepackage{xcolor}
\usepackage{framed}

\spanishdecimal{.}

\setlist[enumerate,1]{label=\textbf{\arabic*.}, leftmargin=*, itemsep=6pt}
\setlist[enumerate,2]{label=(\alph*), itemsep=4pt}
\setlength{\parindent}{0pt}
\setlength{\parskip}{4pt}
\pagestyle{plain}

\newcommand{\Ein}{E_{\text{in}}}
\newcommand{\Eout}{E_{\text{out}}}
\newcommand{\Hyp}{\mathcal{H}}
\newcommand{\dvc}{d_{\text{VC}}}
\newcommand{\vx}{\mathbf{x}}
\newcommand{\vw}{\mathbf{w}}
\newcommand{\vy}{\mathbf{y}}
\newcommand{\R}{\mathbb{R}}
\newcommand{\T}{^{\top}}
\newcommand{\E}{\mathbb{E}}
\newcommand{\Prob}{\mathbb{P}}
\newcommand{\pts}[1]{\hfill\textbf{(#1~puntos)}}
\newcommand{\linea}[1]{\rule{#1}{0.4pt}}
\newcommand{\ptsinc}[1]{{\unskip\nobreak\hfil\penalty50\hskip1em\hbox{}\nobreak\hfill(#1)\parfillskip=0pt\finalhyphendemerits=0\par}}

\newcommand{\encabezado}[2]{%
\begin{center}
{\large\bfseries Reconocimiento de Patrones}\\[2pt]
Universidad de Sonora \textperiodcentered{} Licenciatura en Ciencias de la Computación\\
Semestre 2027-1 \textperiodcentered{} Prof.\ Julio Waissman Vilanova\\[10pt]
{\Large\bfseries #1}%
\ifdefined\tipodoc\\[4pt]{\large\itshape\tipodoc}\fi
\end{center}
\ifdefined\tipodoc\else
\vspace{6pt}
\textbf{Nombre:} \hrulefill\quad \textbf{Expediente:} \rule{3.5cm}{0.4pt}\par
\vspace{6pt}
\fi
#2
\par\medskip\hrule\medskip
}

\newcommand{\espacio}[1]{\par\vspace{#1}}
\newcommand{\nuevapagina}{\newpage}


\begin{document}

\encabezado{Examen rápido 4 \textperiodcentered{} Modelos lineales}{%
Tiempo: 20 minutos. Responde en esta hoja. Puedes usar calculadora; no se permiten apuntes. Deja indicadas las operaciones.}

\begin{enumerate}

%% ---------------------------------------------------------------
\item \textbf{Mínimos cuadrados con cuatro puntos.} Queremos ajustar $h(x) = w_0 + w_1 x$ a los datos
\begin{center}
\begin{tabular}{c|cccc}
$x$ & 0 & 1 & 2 & 3 \\ \hline
$y$ & 1 & 3 & 4 & 8
\end{tabular}
\end{center}
\begin{enumerate}
  \item Escribe la matriz $X$ (con la columna de unos), calcula $X\T X$ y $X\T \vy$, y resuelve las ecuaciones normales.
  \espacio{4.5cm}

  \item Calcula los residuos $r_n = y_n - h(x_n)$ y el error $\Ein = \frac1N \sum_n r_n^2$. Verifica que el vector de residuos es ortogonal a las dos columnas de $X$. ¿Por qué tenía que pasar?
  \espacio{4cm}
\end{enumerate}

\nuevapagina
%% ---------------------------------------------------------------
\item \textbf{Logística y temperatura.} Un modelo (ficticio) de regresión logística estima la probabilidad de que una junta de un cohete sufra daño en función de la temperatura $t$ (en °C) al despegue:
\[
P(\text{daño} \mid t) = \theta(5 - 0.25\, t), \qquad \theta(s) = \frac{1}{1 + e^{-s}} .
\]
\begin{enumerate}
  \item Calcula la probabilidad de daño a $12$\,°C y a $24$\,°C. ¿A qué temperatura vale exactamente $0.5$?
  \espacio{2.8cm}

  \item ¿Por qué factor se multiplican los momios de daño cuando la temperatura \emph{baja} 4\,°C?
  \espacio{2cm}

  \item Con tres vuelos, $(t, y) = (10, 1),\ (20, 0),\ (30, 0)$, y empezando en $\vw = (w_0, w_1) = (0, 0)$, calcula el gradiente de la entropía cruzada, $\nabla\Ein = \frac1N X\T(\mathbf p - \vy)$. ¿En qué sentido moverá a $w_1$ un paso de descenso por gradiente, y por qué tiene sentido?
  \espacio{3.5cm}
\end{enumerate}

%% ---------------------------------------------------------------
\item \textbf{Una regresión de Poisson.} Para el número de bicicletas rentadas en una hora se ajustó $y \mid \vx \sim \operatorname{Poisson}(\mu)$ con
\[
\ln \mu = 2 + 0.05\,\text{temp} - 0.4\,\text{lluvia},
\]
donde temp está en °C y lluvia vale 1 si llueve y 0 si no.
\begin{enumerate}
  \item ¿Cuántas rentas espera el modelo a 20\,°C sin lluvia? ¿Y con lluvia? ¿En qué porcentaje cambian las rentas cuando llueve?
  \espacio{2.5cm}

  \item ¿Qué ventaja tiene este modelo sobre una regresión lineal $\mu = w_0 + w_1\,\text{temp} + w_2\,\text{lluvia}$ para estos datos?
  \espacio{2.2cm}
\end{enumerate}

\end{enumerate}


\end{document}
